\section{Findings}
\label{sec:findings}

Stating the laws and then checking them found six defects.  Each was reproduced
by a failing test before it was fixed.  Only the first is a security defect; the
rest are consistency defects that fail closed or are cosmetic, but each is a
place where two parts of the program disagreed.

\begin{longtable}{@{}p{0.03\linewidth}p{0.30\linewidth}p{0.24\linewidth}p{0.36\linewidth}@{}}
\toprule
& Defect & Law broken & Fix \\
\midrule
\endhead
1 & \textbf{Deny bypass through full case folding.}  On APFS the name
\code{.\ss h} opens \code{.ssh}, but the matcher folded with simple folding, so
a rule for \code{\~{}/.ssh} did not match it; \code{Open}, \code{Head} and the
listing of the parent all served the denied folder.  Also long s, the Kelvin
sign and the ligatures.
& \ref{law:fold-invariance}, \ref{law:absent}
& Rules, roots and search compare after full folding (\code{rules.Fold}); the alias
set was measured on this file system and is a test.  \label{find:fold} \\
2 & \textbf{Listed but unopenable.}  A plain file named like a \code{dir/} deny
rule (a file called \code{cache}) was shown by \code{List} but refused by
\code{Open}, because \code{Open}'s pre-check assumed every path is a directory.
& \ref{law:listed-open}
& The pre-check treats the final component as a file; the precise check after
opening still refuses real directories. \\
3 & \textbf{Glob class matched across a slash.}  \code{[!a]} compiled to
\code{[\^{}a]}, which matches \code{/}, so a pattern with no slash could match
across directory components.
& \ref{law:component}
& Negated classes exclude \code{/}; a \code{/} inside a class is a pattern error. \\
4 & \textbf{Row order not transitive.}  A modified time that could not be parsed
compared equal to everything, so \code{a}$=$\code{b}, \code{a}$=$\code{c} but
\code{b}$<$\code{c}.  The resulting order depended on the input order.
& \ref{law:total}, \ref{law:input-order}
& Unreadable times sort as the oldest; comparisons no longer use subtraction. \\
5 & \textbf{Filter and search disagreed.}  The filter box folded with simple
lower-casing, so \code{strasse} did not find \code{Stra\ss e}, while filename
search did.
& \ref{law:filter}
& The filter folds with lower-, upper-, lower-casing, which matches full folding
on the measured aliases. \\
6 & \textbf{Stale fallback.}  After an image failed to load, the fall-back request
had its own cancellation handle, so a slow answer could replace the preview of a
newer selection.  \label{find:stale}
& \ref{law:fresh}
& One handle per selection is shared by every request started for it. \\
\bottomrule
\end{longtable}

\subsection{Observations that are not defects}

\begin{itemize}
  \item A dangling symbolic link is listed (marked as broken) but cannot be
        opened.  Law~\ref{law:listed-open} excludes it.
  \item Hidden entries are not listed but can be opened, and a hidden directory
        that is entered on purpose shows its contents.
        (Law~\ref{law:hidden-reachable}.)
  \item A symbolic link \emph{named} like a denied file is refused even when its
        target is harmless.  Definition~\ref{def:ok} checks both names.
  \item The macOS by-inode path \code{/.vol/<device>/<inode>} was probed as a
        possible alias that would evade path rules; on the tested system the
        directory is empty and the path cannot be opened.
\end{itemize}

\subsection{What the method did and did not find}

Findings 1 to 3 come from the rule and access laws and are about names: two
spellings, two types, two segments treated as one.  Findings 4 to 6 come from
the frontend laws.  The laws that held on first check (monotonicity, order
independence, ancestor closure, narrowing, search implies listed, denied is
absent, the round trips) are as informative as the failures, since they turn
long-standing assumptions into checked statements.
